%!TEX root = ../thesis.tex
\mychapter{Introduction}{Introduction}{}
\label{chap:introduction}


\mysection{Background}{Background}
\label{sec:background}

A robot that takes spoken instructions removes two obstacles that make most mobile robots
difficult to operate day to day. Typically a mobile robot operates using a screen to navigate, and an operator trained on its command
set. This matters even more for a small indoor rover, because the tasks it can be
asked to do, such as checking whether a door is open or looking at what is on a shelf are exactly the kind of tasks a person would rather say out loud
than type into an application.

\Acp{llm} are neural networks which are based on the transformer architecture, and are trained on very
large collections of text. Through this training the LLM's learns to generate the next word in the
sequence of text given the previous words in the sequence. This makes LLM's well suited to interpreting
everyday instructions as well as have a huge general knowledge base of the world. For example, when given a natural-language command and a description of what a robot can do, an
\ac{llm} can break the command into a sequence of simple actions called a \ac{robotactionprimitive} (RAP). This RAP can be used by any robot, and can be made up of actions that the robot is able to
perform, without a developer having to hand-code a fixed list of recognised
phrases~\cite{ahn2022saycan}. Ambiguous commands can also be resolved through clarifying questions, when a command leaves out information, the model can ask the
user follow-up questions before it commits to a plan. This exchange, in which the user and the
robot take turns to communicate until the task is clear, is referred to in this report as a
\emph{multi-turn dialogue}~\cite{hori2025longhorizon}.

Handing the planning to a language model introduces two problems. The first problem is that spoken commands are
often ambiguous. So systems that always acts on its first interpretation will act on the wrong one often enough to matter~\cite{ren2023knowno}. 
The second problem is that language models hallucinate. They
produce fluent, confident output that does not correspond to anything real, and this becomes
more likely as a task is broken into more steps~\cite{liang2024introspective}. The techniques
that detect ambiguity or hallucination before a plan is set in stone can also be used while the plan is
being carried out, by monitoring the robot's observations for signs that something has gone
wrong~\cite{sinha2024aesop, elhafsi2023semantic}.

The systems that demonstrate these techniques typically rely on large models hosted in the
cloud and accessed over the internet. These systems also typically run on robots with rich sensing such as \acs{lidar},
an \ac{imu} and wheel encoders that measure how far each wheel has turned. While both assumptions are
reasonable in a research laboratory, they do not hold in every setting in which a
voice-commanded robot would be useful.

A use case of running local only models on a robot is in a private home. Every command sent to a cloud-hosted model, and
every camera frame used to check the robot's progress. These commands and images leaves the home and is processed on a
third party's servers, which many households would not accept for a device that moves through
their living space and records it. A home robot also has to be affordable, which rules out
\acs{lidar} and a powerful on-board computer and leaves a camera as the most practical source of
information about its surroundings. Low-cost drive platforms often also lack wheel encoders and
an \ac{imu}, and their wheels slip enough that a command such as ``drive forward one metre''
cannot be trusted to be carried out accurately. In this setting, the language models must run on
hardware the household already owns, such as a desktop computer with a consumer graphics card,
and the robot has to rely on its camera to know whether its plan is working. Where privacy is not
a concern and a reliable internet connection is available, a cloud-hosted model remains the
better choice, since it is larger and more capable than any model that fits on consumer
hardware.

\mysection{Problem Statement and Research Questions}{Problem Statement and Research Questions}
\label{sec:problem-statement}

There is not a full understanding of whether language and vision models small enough to run on a single consumer
desktop computer and can reliably turn a vague spoken command
into a plan that a low-cost indoor rover can carry out using only its camera and simple on-board
sensors. It is also not known how much complexity such a command can contain before the locally
hosted models begin to produce plans that do not match the command, the room or the rover's
capabilities. Existing work either assumes large, cloud-hosted models, or does not connect a
spoken clarification dialogue to a physical robot that must wait for the user's confirmation
before it moves.

In my skripsie, consumer hardware refers to a gaming desktop computer with an AMD Ryzen~5~3600
processor, 16~GB of system memory and an NVIDIA GeForce GTX TITAN~X graphics card with 12~GB of
video memory, representative of a mid-range gaming computer several years old. All locally
hosted models, for speech recognition, planning and vision. This all must run within the memory of this
machine.

To make these questions measurable, the commands the rover must handle are grouped into the
levels of increasing complexity listed in Table~\ref{tab:command-levels}.

\begin{table}[htbp]
    \centering
    \caption{Levels of command complexity considered in this project.}
    \label{tab:command-levels}
    \begin{tabular}{llp{7cm}}
        \toprule
        Level & Type & Example \\
        \midrule
        L1 & Single action & ``Turn left.'' \\
        L2 & Sequence of actions & ``Drive forward, turn left, then drive forward again.'' \\
        L3 & Perception-dependent & ``Drive forward until you see the door.'' \\
        L4 & Conditional & ``Check whether the door is open. If it is, go through it; otherwise, come back.'' \\
        L5 & Vague & ``Go and see who is in the kitchen.'' \\
        \bottomrule
    \end{tabular}
\end{table}

\begin{enumerate}
    \item[\textbf{RQ1}] Can locally hosted models, running on the consumer hardware defined
        above, turn a vague spoken command into a plan that the rover can execute, and which
        levels of command in Table~\ref{tab:command-levels} can they handle reliably?
    \item[\textbf{RQ2}] For each level of command that the models can handle, how many steps can
        be chained into a single command before the plan starts to include hallucinations,
        meaning actions, objects or locations that are not supported by the command, the room
        or the rover's capabilities?
\end{enumerate}

In both questions, the rover's camera is its main source of feedback about the world. A
\ac{vlm} uses camera images to check the plan against the room before the rover departs, and to
confirm the rover's progress while the plan is carried out.

\mysection{Objectives}{Objectives}
\label{sec:objectives}

The main objective of this project is to build and evaluate a fully local pipeline that turns a
spoken, potentially ambiguous command into a plan an indoor rover can carry out, using models
small enough to run on a single consumer-grade machine rather than a cloud API. To reach this,
the project must:

\begin{enumerate}
    \item Assemble a rover from off-the-shelf components on which the rest of the system can be
        tested, with a camera as its main sensor and an ultrasonic distance sensor and line
        sensor in supporting roles.
    \item Determine which speech-to-text models can transcribe spoken commands accurately and
        quickly enough when run on the consumer hardware, including for speakers of South
        African English and Afrikaans.
    \item Implement the rover's capabilities as a set of modular functions, and determine which
        levels of command in Table~\ref{tab:command-levels} a locally hosted planner can
        reliably map onto these functions.
    \item Implement a multi-turn dialogue that resolves ambiguity in a command before a plan is
        generated, followed by a spoken confirmation step in which the rover reads its plan back
        to the user and only starts moving once the user agrees.
    \item Ground each plan in the actual room, using a walkthrough scan recorded beforehand and
        camera frames checked during execution.
    \item Measure and minimise the time taken from the end of a spoken command to a confirmed
        plan, without reducing the correctness of the plan.
    \item Evaluate how the correctness and executability of plans change as commands move up
        the levels in Table~\ref{tab:command-levels} (RQ1).
    \item Evaluate how the number of steps chained in a command relates to the rate of
        hallucination (RQ2).
    \item Compare the locally hosted pipeline against equivalent cloud-hosted models in terms of
        plan correctness and response time.
\end{enumerate}

\mysection{Scope and Limitations}{Scope and Limitations}
\label{sec:scope-and-limitations}

Since this project is concerned with turning vague spoken commands into robot action plans. 
Its scope is limited to the pipeline from speech to executed plan. This includes the transcription of speech, the clarification of ambiguous commands through spoken dialogue, plan generation and verification, and the spoken confirmation step.
 Navigation is used only as a means of carrying out and testing the plans, and improving navigation itself is not an aim of this project.

Before the rover is used in a room, the room is recorded once in a short walkthrough video. A
\ac{vlm} describes the key frames of this video in text, producing a written summary of the objects
in the room and their layout. This summary is given to the planner, so that plans refer only to
things that are actually present. Selected frames are also stored with their descriptions as a
reference set, against which the rover's live camera frames are compared during a mission to
check that it is where the plan expects it to be. The scan is not turned into a
three-dimensional model or map, and building a general-purpose mapping system is outside the
scope of this project.

Cloud-hosted models are used only as a benchmark for the locally hosted pipeline, and are not
part of the system under test.

The rover is deliberately limited to a camera, an ultrasonic distance sensor and a line sensor,
with no wheel encoders, \ac{imu} or \acs{lidar}. This reflects the low-cost home robot described
in Section~\ref{sec:background}, and forces the pipeline to rely on the camera to determine
whether a plan is succeeding, which is the capability this project sets out to test. As a result,
movement is carried out without precise measurement of distance or heading, and commands that
require precise distances beyond what the ultrasonic sensor can measure are not supported.

Similarly, the system is not designed to operate in arbitrary rooms or with arbitrary commands.
Evaluation is limited to a small number of test rooms and a fixed set of test commands at each
level in Table~\ref{tab:command-levels}, which is sufficient to answer RQ1 and RQ2. The results
therefore describe the pipeline's behaviour on the cases tested. Rather than guaranteeing its
performance in other rooms or with other commands.

\mysection{Report Outline}{Report Outline}
\label{sec:report-outline}

The remainder of this report is structured as follows: a review of related work in task
planning, ambiguity resolution, and hallucination in robotic systems; the design of the
speech-to-plan pipeline and its clarification loop; the rover platform and the interface between
the planning software and the hardware it controls; the test methodology, including the
simulated executor used in place of physical hardware; the results and discussion for each
research question; and a concluding summary with recommendations for future work.